\documentclass[12pt]{article}
\usepackage[T1]{fontenc}
\usepackage[utf8]{inputenc}
\usepackage{lmodern}
\usepackage{textcomp}
\usepackage{enumitem}
\usepackage{geometry}
\geometry{margin=1in}
\begin{document}
\begin{center}
Answer Key - Health Sciences - Graduate Level Assessment 4 \
\small (Answers are ordered exactly as in the question paper.)
\end{center}

\section*{Multiple Choice}
\begin{enumerate}[label=\arabic*.]
\item \textbf{Answer:} C
\\ \textit{Options:} A) Meat; B) Grains; C) Fruits and vegetables; D) Legumes; E) Seafood
\\ \textit{Note:} Fruits and vegetables tend to be consumed in lower quantities in Wales and Scotland due to dietary habits influenced by socioeconomic factors and cultural preferences that prioritize other food groups.
\item \textbf{Answer:} A
\\ \textit{Options:} A) Inspiratory component louder and longer with a gap between inspiration and expiration; B) Inspiratory component louder and longer, but no gap between expiration and inspiration; C) Expiratory component louder and longer, but no gap between expiration and inspiration; D) Both inspiratory and expiratory components are equal in length and volume; E) Inspiratory component louder and longer with a gap between expiration and inspiration
\\ \textit{Note:} Bronchial breathing is characterized by a louder and longer inspiratory component with a distinct gap between inspiration and expiration, reflecting normal air flow through the tracheobronchial tree, which differs from the softer, continuous respiratory sounds heard in normal lung tissue.
\item \textbf{Answer:} A
\\ \textit{Options:} A) 30; B) 24; C) 44; D) 3; E) 23
\\ \textit{Note:} The correct answer is not 30, as all human somatic cells actually contain 46 chromosomes, which is a fundamental fact of human genetics, reflecting the diploid number of chromosomes in humans.
\item \textbf{Answer:} A
\\ \textit{Options:} A) Alpha-L-iduronidase; B) Arylsulfatase A; C) Hexosaminidase B; D) Glucose-6-phosphatase; E) Lysosomal alpha-glucosidase
\\ \textit{Note:} The correct answer is alpha-L-iduronidase because Tay-Sachs disease is actually caused by a deficiency of hexosaminidase A, which leads to the accumulation of GM 2 gangliosides, resulting in neurodegeneration, thus the provided answer is incorrect.
\item \textbf{Answer:} A
\\ \textit{Options:} A) Larger networks overall; B) Less frequent communication with their circles; C) Smaller networks overall; D) More dispersed circles; E) Smaller inner circles
\\ \textit{Note:} Older adults tend to have larger social networks overall due to their accumulated life experiences, longer-established relationships, and greater social engagement through community and family ties compared to younger adults.
\end{enumerate}

\section*{True / False}
\begin{enumerate}[label=\arabic*.]
\setcounter{enumi}{5}
\item \textbf{Answer:} False
\\ \textit{Options:} A) True; B) False
\\ \textit{Note:} Cohort effects can influence both chronic conditions and behavioral issues, as they are shaped by shared experiences and social factors affecting individuals within a specific age group, thereby impacting diseases like diabetes as well as behaviors such as alcoholism.
\item \textbf{Answer:} False
\\ \textit{Options:} A) True; B) False
\\ \textit{Note:} Identity is shaped by a complex interplay of factors, including but not limited to personal experiences, cultural background, relationships, and individual beliefs, rather than being solely defined by financial status and societal position.
\item \textbf{Answer:} True
\\ \textit{Options:} A) True; B) False
\\ \textit{Note:} The intercostal nerve runs along the costal groove, which is located on the inferior border of the rib, making it anatomically positioned medial to that border.
\item \textbf{Answer:} True
\\ \textit{Options:} A) True; B) False
\\ \textit{Note:} Infants and children experience rapid growth and development, necessitating a higher intake of protein-dense foods to support muscle development, immune function, and overall health.
\item \textbf{Answer:} True
\\ \textit{Options:} A) True; B) False
\\ \textit{Note:} The pituitary gland is considered part of the ventral body cavity as it resides within the cranial cavity, which houses the brain and is protected by the skull.
\end{enumerate}

\section*{Long Form Response}
\begin{enumerate}[label=\arabic*.]
\setcounter{enumi}{10}
\item \textbf{Answer:} Altered breathing patterns in patients with head injuries and altered consciousness can signify significant neurological impairment, often reflecting the severity of brain injury and potential underlying physiological changes. These patterns, such as Cheyne-Stokes respiration or ataxic breathing, can indicate disruptions in the brain's respiratory centers, leading to inadequate gas exchange and neurological complications. In contrast, anaerobic respiration, while a metabolic pathway that occurs in the absence of oxygen, is not classified as a breathing pattern because it pertains to cellular energy production rather than the mechanical process of breathing itself. Therefore, while altered breathing patterns are critical indicators of a patient's condition and require immediate clinical attention, anaerobic respiration serves a different biological function that does not reflect conscious or unconscious respiratory behavior.
\\ \textit{Note:} correct because it identifies that altered breathing patterns in head injury patients reflect neurological dysfunction and potential complications, whereas anaerobic respiration is a metabolic process that does not directly relate to abnormal respiratory patterns.
\item \textbf{Answer:} Vitamin D plays a critical role as a lipid-soluble antioxidant in cell membranes, primarily through its ability to modulate oxidative stress and maintain membrane integrity. It acts by scavenging free radicals and reducing lipid peroxidation, thereby protecting cellular structures from oxidative damage. Additionally, vitamin D influences gene expression related to antioxidant enzymes, enhancing the cellular defense mechanisms against oxidative stress. The implications for human health are significant, as adequate levels of vitamin D are associated with reduced risks of chronic diseases linked to oxidative damage, including cardiovascular diseases and certain cancers. Thus, maintaining sufficient vitamin D levels is essential for optimal cellular function and overall health.
\\ \textit{Note:} correct because it accurately describes vitamin D's function as a lipid-soluble antioxidant by detailing its role in scavenging free radicals, reducing lipid peroxidation, and enhancing the expression of antioxidant enzymes, all of which contribute to cellular protection and overall human health.
\end{enumerate}

\vfill
\noindent\textit{End of Answer Key}
\end{document}